\begin{abstract}
We reconstruct the technical execution of Iran's 2025--26 nationwide
internet shutdown by jointly measuring BGP routing visibility (RouteViews
and RIS route collectors, 8-hour resolution across a 14-month study period
plus five BGP-update-resolution event windows spanning 2019--2026) and
active-probing/network-telescope outage signals (IODA), aligned per
Autonomous System. Every analytical threshold, exclusion rule, and phase
boundary used is data-derived and recorded in an append-only decision log
before its effect on any result is known, including two corrections
surfaced by that discipline during this study: a double-counting bug in
address-space accounting for overlapping BGP announcements, and a dark-
classification threshold that had been left as an unfinalized placeholder
past the point its own governing rule required closing it. We find the
blackout was filtering-dominant: routes remained visible in BGP while
reachability collapsed, not a control-plane withdrawal -- a pattern absent
from a matched control population of three regional neighbors and
independently corroborated by a contemporaneous study using an unrelated
measurement stack. Restoration during the subsequent partial, tiered
recovery phase was selective by sector, though coarser than a clean
per-type ranking: government and core state infrastructure recovered
fastest and most completely, an unclassified long tail slowest and least
completely, with genuine complications reported openly, not smoothed over.
We report disconfirming evidence and open limitations explicitly, including
an untested topological sub-claim and one comparison lacking a
control-population check, and tier every empirical claim by the strength of
evidence it rests on.
\end{abstract}
